\section{Conclusion}
\label{sec:conclusion}

We presented MultiEMiT, a multimodal dataset of Italian Twitter commentary annotated across five pragmatic layers: Plutchik emotions, irony, offensiveness, target, and semiotic relations. Benchmarking specialized Dual-Encoders against Pixtral 12B revealed three key findings:
\begin{itemize}
    \item \textbf{RQ1 (Visual Context):} Imagery resolves textual ambiguity when modalities align, substantially improving multi-label emotion recognition over blind text baselines. In irony detection, however, unguided visual features introduce semantic noise and degrade performance relative to text alone due to the modality gap.
    \item \textbf{RQ2 (Correlation Prompting):} Injecting empirical correlation priors as textual prefixes consistently improves classification across tasks. For irony, correlation prompting resolves cross-modal conflict and lifts the Dual-Encoder to the highest overall performance.
    \item \textbf{RQ3 (Architectural Paradigms):} Compact discriminative models (BERT-ViT with Classifier Chains) outperform the 12B foundation model on domain commentary. Pixtral exhibits exposure bias on long-tail classes, requiring post-hoc truncation ($n=2$) to curb hallucinations.
\end{itemize}

Looking ahead, our findings open two avenues for future work. First, expanding MultiEMiT beyond political talk shows (\textit{Propaganda Live}) and reality television (\textit{Pechino Express}) to scripted drama, live sports, and cultural broadcasts would help reduce genre-specific class skew and test cross-genre generalization. Second, future studies could investigate affective annotation schemas designed specifically for Social TV dynamics. Rather than relying solely on general categorical taxonomies like Plutchik's, dedicated schemas could target engagement-driven states (\textit{hype}, \textit{amusement}, \textit{frustration}, \textit{empathy}, \textit{boredom}) alongside viewer stances (\textit{support}, \textit{opposition}, \textit{hatewatching}, \textit{neutral}) to better capture how digital television audiences react.
